\chapter{State of the Domain}

\section{Background on Fraud Detection and Imbalanced Domains}
Modern fraud detection systems are designed to distinguish legitimate customer behaviour from fraudulent activity in near real time. Practitioner literature emphasises telemetry-rich features, cascading rules layered with learner scores, and tight coupling to investigation workflows \cite{fraud_handbook}. Survey syntheses specialised to payment-card contexts stress that methodological diversity (from shallow scoring models to recurrent or graph-augmented representations) often matters less than consistent treatment of labels, leakage, and deployment skew—because reported offline gains rarely translate when the fraud process itself shifts under detection pressure \cite{lucas_jurgovsky_ccfraud_survey2020}.

The dominant statistical difficulty remains extreme imbalance: positives are so sparse that empirical risk minimisers aligned with misclassification counts learn to mimic the majority label. A classifier that abstains from ever raising alerts can exhibit deceptively high accuracy while economically failing the institution \cite{he_garcia_imbalanced2009,fraud_handbook}. Formal surveys of imbalance learning group countermeasures into \emph{data-level} approaches (resampling, informed synthetic minority generation), \emph{algorithm-level} adaptations (modified kernel methods, ensemble strategies, cost-sensitive objectives), and hybrid combinations; the guiding principle is to prevent optimisation from treating rare positives as ignorable noise \cite{he_garcia_imbalanced2009}. In fraud, these corrections interact with temporal structure: stratified holdout along time and merchant segments is typically as important as the choice between, say, oversampling versus class-weighted loss, because naive shuffles inflate precision by letting models memorise ephemeral cohort effects \cite{lucas_jurgovsky_ccfraud_survey2020,fraud_handbook}.

Robust pipelines therefore pair domain-aware feature construction—amount, velocity, merchant category, geo-temporal signals, device identity proxies—with metrics that make the minority class visible in evaluation rather than hiding it inside aggregate accuracy. Average precision under the precision--recall curve is one standard scalar; its interpretive role is tied to how analysts spend investigation budget at different score cutoffs, not to a platonic ``best'' model ordering independent of operating costs \cite{davis_goadrich_icml06,fraud_handbook}.

\section{Evolution of Machine Learning Approaches}
Historical payment fraud stacks layered deterministic rules atop relational databases—velocity caps, geographically implausible hops, deterministic blocklists—with human-authored thresholds tuned reactively following loss events \cite{fraud_handbook}. The fragility manifests when adversaries diversify patterns faster than committees encode exceptions; escalating rule counts inflate maintenance costs while degrading investigative signal-to-noise ratios.

\textbf{Phase shift toward supervised surrogates.} By the middle of last decade empirical papers increasingly trained supervised classifiers directly on adjudicated transactional archives \cite{lucas_jurgovsky_ccfraud_survey2020}. Random forests, gradient-boosted decision trees (GBDT families), elastic-net regularised logistic models, kernel methods, and—more recently—shallow neural networks competed on ROC or PR summaries under severe imbalance regimes. Institutional adoption aligned with toolchain maturity (.NET-centric, JVM, and scientific Python ecosystems) rather than ideological purity concerning interpretability versus accuracy.

\textbf{Complementary anomaly paradigms.} Label scarcity for novel attacks motivates unsupervised anomaly scores (isolation forests, Gaussian mixtures, spectral techniques) calibrated with extreme-value heuristics. Hybrid deployments treat anomaly triggers as widening candidate pools upstream of discriminative scorers, reflecting operational reality that pure supervision lags zero-day fraud until labels arrive \cite{yousefi_fraud_survey2019,fraud_handbook}.

\textbf{Representation learning frontiers.} Graph neural models encode merchant--card bipartite structures; recurrent and transformer-based encoders consume event sequences; contrastive embeddings cluster device fingerprints. These directions expand hypothesis spaces but escalate explainability burdens because attributions diffuse across latent states \cite{lucas_jurgovsky_ccfraud_survey2020}. The dissertation references them sparingly whenever narrative clarity risks dilution—the focal predictors remain tabular stochastic maps consistent with PCA-compressed benchmarking corpora.

Despite architectural diversity, nonlinear predictors converge on opaque decision boundaries unless structural constraints or ancillary explanation layers externalise curvature. Subsequent sections formalise taxonomy so readers can reconcile competing interpretability philosophies without conflating them.

\section{Semantic Taxonomy of Model Interpretability}
Explainability notions are heterogeneous even before one commits to fraud-specific jargon; Guidotti \emph{et al.}\ organise the space by aligning \emph{(i)} problem types (inspect global behaviour versus justify a singleton decision), \emph{(ii)} black-box classes, and \emph{(iii)} target explanation formats—stressing that each algorithmic proposal embeds tacit semantic commitments about what counts as understandable \cite{guidotti_blackbox_survey2018}. The present thesis adopts terminology consistent with overlapping survey traditions \cite{xai_survey,molnarInterpretableMachineLearningBook}.
\begin{enumerate}
    \item \textbf{Transparency depth:} \emph{ante-hoc} / intrinsic interpretability (models whose structure directly supports auditing) versus \emph{post-hoc} interpretation (explainers applied after training).
    \item \textbf{Scope:} \emph{global} understanding summarises unconditional dependencies between predictors and risk; \emph{local} understanding conditions on a realised transaction $x$ and investigates sensitivity of $\hat{y}$ under controlled perturbations.
    \item \textbf{Audience:} compliance-oriented narratives differ from fidelity-oriented diagnostics demanded by modelling teams; ambiguity between \emph{plausible} summaries and \emph{faithful} reconstructions motivates explicit caveats whenever explanations feed operational decisions \cite{xai_survey}.
\end{enumerate}
Fraud operations primarily seek local, post-hoc evidence per alert; therefore later chapters analyse methods that quantify instance-level contributions while referencing global artefacts only as corroborative context.

\section{The Paradigm of Post-Hoc Analysis}
Among post-hoc families, perturbation-derived surrogates and axiomatic attribution maps dominate tabular deployments. Conceptually these methods analyse the directional derivatives of scoring maps $f_\theta:\mathbb{R}^d\to[0,1]$ induced by stochastic classifiers fitted on historical fraud labels.

These techniques treat---at least nominally---the underlying predictor as inaccessible except through query-oracle evaluations. Inputs are manipulated, neighbourhoods are calibrated, cooperative games over feature coalitions may be enumerated (or approximated), and explanatory vectors summarise emergent curvature of $f_\theta$ restricted to neighbourhoods of empirical interest.

The methodological advantage motivating post-hoc choice in financial tooling is encapsulation: the interpretability façade can nominally traverse model upgrades conditioned on preserving API contracts---subject to explanatory stability constraints discussed later.

\section{Advanced Post-Hoc Interpretability Frameworks}
Two of the most prominent model-agnostic, post-hoc explanation frameworks currently driving the field are LIME and SHAP. Both frameworks can be applied to a broad range of predictive models and are widely adopted in explainable AI practice \cite{lime_original,shap_original,xai_survey}.

\subsection{Local Model-Agnostic Explanations (LIME)}
LIME is a post-hoc technique that approximates the predictions of any underlying black-box model by training an interpretable model locally around the prediction of interest. The motivating smoothness heuristic is heuristic rather than causal: nonlinear boundaries may appear roughly linear inside small balls in feature space, but the radius and shape of those balls materially determine which surrogate coefficients analysts see \cite{lime_original,guidotti_blackbox_survey2018}.

LIME perturbs the input data points to generate a localized dataset, records the black-box model's responses to these variations, and builds a weighted linear model based on the proximity to the original instance. The explanation is generated by minimizing the following objective function:
\begin{equation}
\xi(x) = \underset{g \in G}{\operatorname{argmin}} \mathcal{L}(f, g, \pi_x) + \Omega(g)
\end{equation}
Where $\mathcal{L}$ is a loss function measuring how unfaithful the surrogate model $g$ is in approximating the black box $f$ in the locality defined by $\pi_x$, and $\Omega(g)$ is a complexity penalty ensuring $g$ remains interpretable to human investigators. When tabular perturbations violate feature dependence structure (e.g., independently scrambling amount and merchant ID), surrogates may optimise mathematical fidelity to an oracle evaluated on unrealistic counterfactuals; fraud teams mitigate this pragmatically via background sampling from reference baskets of transactions, explicit rule gates, or hybrid checks against global partial dependence—none of which remove the approximation problem altogether \cite{lime_original,guidotti_blackbox_survey2018}.

\subsection{SHapley Additive exPlanations (SHAP)}
While LIME focuses on local approximation, SHAP provides a unified measure of feature importance based on cooperative game theory. It treats the features of a transaction as "players" in a game where the "payout" is the final fraud probability score. 

SHAP calculates the exact marginal contribution of each feature across all possible feature combinations. The Shapley value $\phi_i$ for a feature $i$ is mathematically defined as:
\begin{equation}
\phi_i = \sum_{S \subseteq N \setminus \{i\}} \frac{|S|! (|N| - |S| - 1)!}{|N|!} \left[ f_x(S \cup \{i\}) - f_x(S) \right]
\end{equation}
Where $N$ is the set of all features, and $S$ is a subset of features not containing $i$. This guarantees that the sum of the individual feature attributions perfectly equals the difference between the model's base expectation and the final prediction—\emph{if} coalition values stem from coherent choices of marginalisation across missing coordinates and baselines.

Exact enumeration is exponential in $|N|$; practical SHAP tooling for tree ensembles approximates cooperative games via polynomial-time procedures at the expense of strictly weaker guarantees outside the asymptotic regimes where sampling estimators concentrate \cite{shap_original}. Analyst-facing dashboards should therefore treat SHAP bars as stable only when score functions vary smoothly and background distributions track live traffic; under sharp concept drift, yesterday's baseline expectation is a moving target, and attributions inherit that movement \cite{lucas_jurgovsky_ccfraud_survey2020,xai_survey}. Adversarial scaffolding further shows that classifiers can retain biased behaviour while LIME/SHAP narratives appear innocuous, underscoring the need for stress tests beyond single-instance reviews \cite{slack2020_fooling_lime_shap}. Operational teams therefore pair perturbation-derived attributions with documentation norms such as Model Cards/Datasheets for datasets when summarising lineage and limitations \cite{mitchell2019model_cards,gebru2021datasheets}.

Outside controlled research settings, interpreting ``good'' explanations requires articulated evaluation regimes (who benefits, what error is tolerated, when recourse is required) rather than informal satisfaction \cite{doshi_velez_kim_interpretability_2017}. Fraud analytics inherits that obligation when investigators rely on ranked features for queue prioritisation.

\section{Measuring Classification Quality under Rare Positives}
Class prevalence $\pi=P(y{=}1)$ in retail card streams is microscopic; thus the contingency matrix induces skewed ROC dominance without exposing precision collapse. Accuracy $\frac{TN+TP}{N}$ inflates mechanically when negatives dominate even if $TP\to0$.

The precision--recall curve traces $\text{Recall}(\tau)$ against $\text{Precision}(\tau)$ as the score threshold varies. Integrated measures such as average precision quantify expected precision across recall operating points without assuming symmetry of errors:
\begin{equation}
\text{AP} = \sum_{k} (R_k-R_{k-1}) P_k ,
\end{equation}
where ordered thresholds induce precision $P_k$ at recall breakpoints $R_k$---a pragmatic scalar aligned with asymmetric costs \cite{fraud_handbook}. ROC-AUC still measures ranking separability among positives versus negatives yet can remain high despite economically unacceptable false-positive tides; therefore both ROC and precision--recall summaries deserve parallel reporting.

Operational decisions hinge on thresholds $\tau$ mapping scores to escalate/stop workflows. Equivalent viewpoints cast $\tau$-selection as constrained optimisation balancing expected fraud loss versus investigation budget: denote by $c_{FP}$, $c_{FN}$ asymmetric costs attached to erroneous legitimate vs fraudulent classifications; heuristic Bayes-risk style criteria motivate studying iso-cost surfaces in ROC or PR geometries even when plug-in probabilities are imperfectly calibrated \cite{fraud_handbook}.

\subsection{Geometry of ROC versus precision--recall under skew}
Davis and Goadrich derive the structural link between ROC curves and precision--recall curves for a fixed scoring classifier: each curve is a monotonic transform of the other through the underlying contingency counts, yet their visual emphases diverge when negatives dwarf positives \cite{davis_goadrich_icml06}. For fraud monitoring, the practical import is twofold. First, PR curves foreground the behaviour of precision at low false-alarm budgets—precisely the region where investigation teams live. Second, dominance arguments that feel intuitive in ROC space (``model A is uniformly better'') do not automatically carry identical operational meaning when translated into PR coordinates; reporting both families of curves mitigates premature model selection based on a single geometric lens \cite{davis_goadrich_icml06,fraud_handbook}.

\subsection{Probability statements and calibration}
Many stacks export $\hat{y}\in[0,1]$ as ``probability'' even when training loss does not enforce reliability; misranking is then conflated with miscounting expected fraud among transactions above a cutoff. Reliability diagrams, Platt scaling, isotonic regression, and related post-training adjustments have long been studied as ways to align predicted frequencies with empirical frequencies on held-out validation sets \cite{niculescu_caruana_calib2005,molnarInterpretableMachineLearningBook}. Calibration quality matters coherently with explanation layers: if $\hat{y}$ is only ordinal, attributions may still rank features, but threshold narratives phrased in percentage language become legally and semantically fragile.

\section{Faithfulness, Stability, and Adversarial Drift}
Local explanations hinge on neighbourhoods; minor distributional perturbations---including adversarial adaptations by fraud actors---alter attributions unpredictably unless regularity conditions (Lipschitz moderation, calibrated scores) partially hold \cite{xai_survey}. Therefore explanations must never be anthropomorphised as causal attributions independent of latent confounders in observational transactional data.

From a longitudinal perspective concept drift induces time-varying $f_{\theta}$, rendering static explanation templates stale. Survey-level accounts of card-fraud ML emphasise drift and adversarial co-evolution as first-class risks rather than edge cases: once a pattern enters training corpora, perpetrators rotate tactics, invalidating both score distributions and the conditional structure on which LIME-like neighbourhoods were calibrated \cite{lucas_jurgovsky_ccfraud_survey2020,xai_survey}. Periodic recalibration, monitoring of covariate shift metrics, joint review of ROC/PR operating points alongside explanation stability dashboards, and revalidation sessions with investigative stakeholders are methodological obligations orthogonal to intermittent accuracy bumps on static benchmarks.

\section{Explainability in Regulated Financial Contexts}
Financial fraud detection is not only a predictive analytics problem, but also a governance problem. In production settings, model outputs may influence high-impact decisions such as transaction blocking, temporary account restrictions, and escalation to manual investigations. As a consequence, institutions must balance statistical performance with accountability, traceability, and decision transparency.

The General Data Protection Regulation (GDPR) and related governance frameworks motivate stronger controls around automated decision support, especially when decisions can materially affect users \cite{gdpr_text}. In practical terms, this means fraud detection platforms should retain model versions, preserve decision logs, and provide interpretable justification artifacts that can be reviewed by analysts and auditors.

\section{Research Gap and Technical Implications}
The literature shows consistent progress in high-performing fraud classifiers, but there remains a practical gap between offline predictive performance and operationally usable explainability. Three technical tensions are especially relevant for this dissertation:
\begin{enumerate}
    \item \textbf{Imbalance versus interpretability:} highly optimized models can improve recall but may produce explanations that are difficult for non-technical investigators to interpret.
    \item \textbf{Local fidelity versus stability:} local explanation methods can be informative for individual cases, yet explanation outputs may vary under small perturbations of the input distribution.
    \item \textbf{Model quality versus workflow cost:} increasing detection sensitivity can raise false-positive volume, directly affecting analyst workload and review latency.
\end{enumerate}

Addressing these tensions requires a design that treats explainability as an integral system requirement rather than a post-deployment add-on. This perspective motivates the architecture proposed in the following chapter, where prediction, explanation, and analyst interaction are developed as a unified workflow.

\section{Institutional uptake patterns and methodological heterogeneity}
Regulated lenders rarely publish detailed detection stacks; fragmented conference literature therefore skews toward academic corpora versus proprietary telemetry. Observational contrasts nonetheless reveal recurring motifs: multidisciplinary squads bridging risk management, AML compliance, cybersecurity, privacy counsel, data engineering, statistics, and behavioural science collaboratively govern model inventories \cite{fraud_handbook}. Explainability artefacts surface variably—as inline SHAP overlays in Jupyter notebooks scrutinised quarterly, textual narratives bundled with challenger models in credit committees, or interactive dashboards tethered to case-management tooling.

\textbf{Operational maturity gradients.} Smaller PSPs prioritise turnkey cloud APIs with interpretability toggles mandated by contractual SLAs; global networks invest in heterogeneous cascades marrying rules, ensembles, anomaly scores, manual overrides, graph analytics, yet share common pain points around alert fatigue and explanation fatigue \cite{yousefi_fraud_survey2019}. Scholars comparing algorithmic manuscripts must therefore relativise leaderboard metrics against institutional envelopes—datasets compress complexity; deployment multiplies interfaces.

\textbf{Evidence gaps.} Transparent replication packages remain rare owing to secrecy compulsion; anonymity trades away semantic variables for reproducibility paradoxically limiting external validity critiques \cite{lucas_jurgovsky_ccfraud_survey2020}. Dissertations like the present one navigate that compromise by marrying public benchmarks with explicit caveats whilst supplying reference software showing how methodological obligations might appear if telemetry constraints relaxed.

Bridging heterogeneous practices, the ensuing chapters articulate normative methodological commitments—threshold literacy, explanatory labelling discipline, versioning transparency—whose satisfaction can be audited even when granular industrial metrics remain unavailable.

\section{Sequential decision processes and cascading policies}
Operational fraud mitigation rarely terminates at a solitary classifier invocation. Telemetry frequently traverses cascading filters: cryptographic authentication, velocity rules, anomaly triggers, blacklist membership, behavioural biometrics scores, culminating in probabilistic escalation engines orchestrated inside case-management dashboards \cite{fraud_handbook}. Each stage transforms the feature manifold available to downstream learners; mistakenly attributing explanatory credit to terminal models alone overlooks conditioning introduced by predecessors.

Sequential testing also imposes multiple-comparison biases when analysts iteratively tweak rules based on exploratory dashboards—a phenomenon analogous to selective inference in biomedical trials. Documenting exploratory versus confirmatory analytic phases aligns dissertations that otherwise risk overstating incidental improvements discovered through repeated leaderboard queries \cite{lucas_jurgovsky_ccfraud_survey2020}.

\textbf{Delayed labels and weak supervision.} Fraud labels finalize asynchronously after investigative cycles; semi-supervised and positive-unlabelled regimes therefore intrude even when researchers idealise fully supervised archives \cite{yousefi_fraud_survey2019}. Methodological chapters cannot ignore these frictions without inviting reviewer scepticism about optimistic PR-AUC claims inferred from provisional positives.

\textbf{Human-in-the-loop adaptation.} Investigators supply feedback tags that reflow into retraining corpora, closing a feedback loop between model outputs and future labels. Game-theoretic perspectives highlight how fraud actors react to deployed policies, invalidating i.i.d.\ sampling heuristics assumed in textbook evaluation \cite{lucas_jurgovsky_ccfraud_survey2020}. Dissertations emphasising static benchmarks should foreground this limitation rather than append it as an afterthought.

These institutional realities motivate the hybrid methodology proposed later: mathematically precise treatment of tabular classifiers coupled with explicit caveats about upstream cascades, delayed supervision, and adversarial adaptation absent from anonymised corpora.